\documentclass{article}
\usepackage{decomp}
\begin{document}
\sourcefile{pesce2022\_supplement.pdf}
\decompnote{Supplementary Information (publisher file 42003\_2022\_3390\_MOESM2\_ESM.pdf), 17 pages, the authors' Word document: a title page, Supplementary Figures 1--14 (one per page, each an image with vector plots where there are plots, legend below) and Supplementary Tables 1--2. The running header ``Pesce et al.'' and the page number at the foot of every page are left out. Text printed inside the figures (panel labels, axis labels, legends) is described in the \alttext{} blocks; the printed figure titles of Supplementary Figs. 8 and 9 and the value table inside Supplementary Fig. 13 are kept as text.}
\pdfpage{1}
\section*{Supplementary Information}

\begin{center}\textbf{3D molecular phenotyping of cleared human brain tissues with light-sheet fluorescence microscopy}\end{center}

\begin{center}\textbf{Authors}\end{center}

\begin{center}Luca Pesce\textsuperscript{1,2}, Marina Scardigli\textsuperscript{1,2}, Vladislav Gavryusev\textsuperscript{1,2}, Annunziatina Laurino\textsuperscript{1,7}, Giacomo Mazzamuto\textsuperscript{1,3}, Niamh Brady\textsuperscript{1}, Giuseppe Sancataldo\textsuperscript{1}, Ludovico Silvestri\textsuperscript{1-3}, Christophe Destrieux\textsuperscript{4}, Patrick R. Hof\textsuperscript{5}, Irene Costantini*\textsuperscript{1,3,6}, Francesco S. Pavone\textsuperscript{1,2,3}\end{center}

\pdfpage{2}
\begin{center}
\includegraphics[width=0.95\linewidth,height=0.6\textheight,keepaspectratio]{figures/sfig1.png}
\end{center}
\noindent \textbf{Supplementary Figure 1. Characterization of SWITCH processed human brain slices treated with different autofluorescence elimination reagents.} Columns 1-4 correspond to the 4 different excitation lights; column 5 correspond to the merge channels for each treatment. The look-up tables (LUTs) are fixed for each treatment. Scale bar = 100 μm.
\alttext{Supplementary Figure 1 (p.\,2) is a grid of 65 small LSFM images, 13 rows by 5 columns. Column headers: 405, 488, 561, 638, Merged channels. Row labels (bold, at the left): SHIELD; SWITCH; SWITCH-0.6\% H\textsubscript{2}O\textsubscript{2}; SWITCH-15\% H\textsubscript{2}O\textsubscript{2}; SWITCH-30\% H\textsubscript{2}O\textsubscript{2}; SWITCH-AR; SWITCH-25\% Quadrol; SWITCH-6 mM Ascorbic acid; SWITCH-20 mM Hydroquinone; SWITCH-SB; SWITCH-NaNH\textsubscript{4} [sic, for NaBH\textsubscript{4}]; SWITCH-1 mM CuSO\textsubscript{4}; SWITCH-10 mM CuSO\textsubscript{4}. Each image is a horizontal strip of the slice cross-section (the illuminated edge at the left) in a single-colour look-up table: blue (405), green (488), yellow (561), red (638); the merged column overlays them (mostly yellow-green). The 405 images are nearly black with a faint blue edge; the 488 column is brightest for SWITCH, SWITCH-AR, ascorbic acid and hydroquinone (bright green over the left half) and dimmest for SHIELD, SB, NaBH\textsubscript{4} and 10 mM CuSO\textsubscript{4}; the 561 and 638 images show scattered bright puncta (lipofuscin granules) on a diffuse background, strongest in the SWITCH, AR, ascorbic acid and hydroquinone rows and weakest in the SB and CuSO\textsubscript{4} rows. White scale bar (100 µm) in the top-left image. No quantitative axes.}

\pdfpage{3}
\begin{center}
\includegraphics[width=0.95\linewidth,height=0.6\textheight,keepaspectratio]{figures/sfig2.png}
\end{center}
\noindent \textbf{Supplementary Figure 2. Effect of tissue processing (SWITCH, SHIELD, and SHORT), autofluorescence treatments (CuSO}\textsubscript{4}\textbf{), and buffer on NeuN fluorescence labeling.} \textbf{a} High resolution imaging of NeuN labeling of different experiments shown in Fig. 1. Scale bar = 100 μm. \textbf{b} Profile plots of NeuN labeling of SHIELD-processed slices incubated with BSA + PBST, treated with AR or H\textsubscript{2}O\textsubscript{2}. (\textbf{c}) Comparison of SWITCH-processed slices treated with H\textsubscript{2}O\textsubscript{2}, AR or by combining both (SHORT). \textbf{d} Profile plots of NeuN antibody incubated with HEPES buffer, or HEPES buffer supplemented with 2.5\% Quadrol and 0.5 M urea in SWITCH/SHIELD-processed slices. \textbf{e} Profile plots of SWITCH-processed slices treated with 1 mM CuSO\textsubscript{4}after NeuN labeling (black and red profiles) or before NeuN labeling (blue and green profiles). Antibodies dilution for all experiments: NeuN 1:100; Alexa Fluor 647 1:500. Incubation temperature of 37 °C (N = 3). Excitation light, 638 nm at a 5 mW laser power.
\alttext{Supplementary Figure 2 (p.\,3): panel a is a raster image column, panels b--e are vector plots. (a) Headed ``NeuN labeling'' with a black bar above, eight horizontal strips of the slice cross-section in a red-yellow (``hot'') look-up table, each labelled at its top right: SHORT – 37 °C (bright yellow cells over the whole width), SHORT – 4 °C (a bright yellow band at the left edge, dark beyond), SHORT – RT (bright cells over most of the width), SWITCH – 37 °C (dim), SHIELD – 37 °C (bright edge only), SHIELD – AR (dim with few cells), SHIELD – H\textsubscript{2}O\textsubscript{2} (dim red with cells), SHIELD – BSA (bright yellow cells throughout); white scale bar (100 µm) in the first strip. (b)--(e) Line plots, x axis ``x position (µm)'' 0--700, y axis ``Intensity'', each series a mean line with a translucent ±SD band and a legend at the upper right. (b) y 0 to 7x10\textsuperscript{4}: ``SHIELD – AR'' (black) ≈5x10\textsuperscript{3} at ≈50 µm, below 2x10\textsuperscript{3} after 150 µm; ``SHIELD – H\textsubscript{2}O\textsubscript{2}'' (red) ≈7x10\textsuperscript{3} at 30--60 µm, decaying to ≈1x10\textsuperscript{3} by 250 µm; ``SHIELD – BSA + PBST'' (blue, cyan band) a tall peak ≈6.5x10\textsuperscript{4} at ≈30 µm, ≈1x10\textsuperscript{4} at 100 µm, ≈2x10\textsuperscript{3} by 200 µm and near zero beyond 400 µm. (c) y 0 to 1,0x10\textsuperscript{4}: ``SWITCH – H\textsubscript{2}O\textsubscript{2}'' (black) ≈5x10\textsuperscript{3} at 50--100 µm, falling to ≈1x10\textsuperscript{3} at 500 µm and ≈5x10\textsuperscript{2} at 700 µm; ``SWITCH – AR'' (red) ≈3x10\textsuperscript{3} at ≈30 µm then ≈2x10\textsuperscript{3}, down to ≈5x10\textsuperscript{2}; ``SHORT – 37 °C'' (blue, cyan band) ≈4x10\textsuperscript{3} at ≈30 µm, ≈3x10\textsuperscript{3} at 100--200 µm, ≈2x10\textsuperscript{3} at 300--500 µm, ≈1x10\textsuperscript{3} at 700 µm. (d) y 0 to 7x10\textsuperscript{4}: ``SWITCH – Hepes buffer'' (black) ≈1x10\textsuperscript{4} at ≈100 µm, decaying, with a second peak ≈1.3x10\textsuperscript{4} at ≈650 µm; ``SWITCH – Hepes buffer + Q/U'' (red) ≈2x10\textsuperscript{4} at ≈50 µm with a wide pink band, ≈1x10\textsuperscript{4} at 100--150 µm, and ≈1x10\textsuperscript{4} at ≈650 µm; ``SHIELD – Hepes buffer'' (blue) ≈1.7x10\textsuperscript{4} at ≈80 µm, falling to ≈2x10\textsuperscript{3} by 200 µm; ``SHIELD – Hepes buffer + Q/U'' (green) a peak ≈5x10\textsuperscript{4} at ≈80 µm with a wide green band reaching 7x10\textsuperscript{4}, ≈5x10\textsuperscript{3} by 200 µm, and a secondary peak ≈1x10\textsuperscript{4} at ≈580 µm. (e) y 0 to 1,0x10\textsuperscript{4}: ``SWITCH – 1 mM CuSO\textsubscript{4} Post-Label'' (black) ≈4x10\textsuperscript{3} at 30--60 µm, ≈5x10\textsuperscript{2} by 300 µm; ``SWITCH – 10 mM CuSO\textsubscript{4} Post-Label'' (red) ≈1.5x10\textsuperscript{3} at ≈30 µm then ≈5x10\textsuperscript{2}; ``SWITCH – 1 mM CuSO\textsubscript{4}'' (blue) ≈1.5x10\textsuperscript{3} at ≈30 µm, ≈1x10\textsuperscript{3}, rising again to ≈1.5x10\textsuperscript{3} at ≈600 µm; ``SWITCH – 10 mM CuSO\textsubscript{4}'' (green) ≈5.5x10\textsuperscript{3} at 40--60 µm, decaying to ≈5x10\textsuperscript{2} by 300 µm.}
\figuredata{b}{figures/sfig2b.csv}{figures/sfig2b.py}{vector paths in the PDF (tier 2), exact up to the PDF's precision: one tiny marker per data point gives the mean, the two polylines of the series colour give the ±SD band edges, axes calibrated on the labelled tick marks (residuals, band-symmetry check and uncertainty in the CSV header)}
\figuredata{c}{figures/sfig2c.csv}{figures/sfig2c.py}{vector paths in the PDF (tier 2), exact up to the PDF's precision: one tiny marker per data point gives the mean, the two polylines of the series colour give the ±SD band edges, axes calibrated on the labelled tick marks (residuals, band-symmetry check and uncertainty in the CSV header)}
\figuredata{d}{figures/sfig2d.csv}{figures/sfig2d.py}{vector paths in the PDF (tier 2), exact up to the PDF's precision: one tiny marker per data point gives the mean, the two polylines of the series colour give the ±SD band edges, axes calibrated on the labelled tick marks (residuals, band-symmetry check and uncertainty in the CSV header)}
\figuredata{e}{figures/sfig2e.csv}{figures/sfig2e.py}{vector paths in the PDF (tier 2), exact up to the PDF's precision: one tiny marker per data point gives the mean, the two polylines of the series colour give the ±SD band edges, axes calibrated on the labelled tick marks (residuals, band-symmetry check and uncertainty in the CSV header)}

\pdfpage{4}
\begin{center}
\includegraphics[width=0.95\linewidth,height=0.6\textheight,keepaspectratio]{figures/sfig3.png}
\end{center}
\noindent \textbf{Supplementary Figure 3. Antibody optimization of SHORT and SHIELD-}\textbf{processed slices. a} High resolution images acquired by LSFM of SHORT and SHIELD-processed slices, labeled for NeuN with Alexa Fluor 647 at different dilution (1:50 and 1:100 for NeuN; from 1:200 to 1:1000 for Alexa Fluor 647). Scale bar = 100 μm. \textbf{b} Profile plot along depth of the two best conditions of NeuN and Alexa Fluor 647 antibodies in SHORT-processed tissues. \textbf{c}-\textbf{d} Profile plots of different antibody dilutions of SHORT-processed slices. \textbf{e}-\textbf{f} Profile plot of different antibody dilutions of SHIELD-processed slices. Incubation temperature of 37 °C (n = 3). Excitation light, 638 nm; laser power, 5 mW.
\alttext{Supplementary Figure 3 (p.\,4): panel a is a raster image block, b--f are vector plots. (a) Two columns of six red LSFM strips, headed ``SHORT'' and ``SHIELD/TDE'' (separated by a blue vertical line), each strip labelled at its top left: NeuN 1:50; SA 1:200 / NeuN 1:50; SA 1:500 / NeuN 1:50; SA 1:1000 / NeuN 1:100; SA 1:200 / NeuN 1:100; SA 1:500 / NeuN 1:100; SA 1:1000 (SA = secondary antibody). The SHORT strips show red cells across the width (brightest for 1:50; SA 1:200 and 1:100; SA 1:200), the SHIELD/TDE strips a bright red band at the left edge with little signal beyond; white scale bar (100 µm) in the first strip. (b)--(f) Line plots, x ``x position (µm)'' 0--700, y ``Intensity'', mean line with ±SD band, legend at the upper right. (b) y 0 to 6,0x10\textsuperscript{3} (ticks 2,0x10\textsuperscript{3}, 4,0x10\textsuperscript{3}, 6,0x10\textsuperscript{3}): ``SHORT NeuN 1:100; SA 1:500'' (black) ≈4.5x10\textsuperscript{3} at ≈40 µm, ≈3x10\textsuperscript{3} at 100 µm, ≈2x10\textsuperscript{3} at 300 µm, ≈1.5x10\textsuperscript{3} at 500 µm, ≈5x10\textsuperscript{2} at 700 µm; ``SHORT NeuN 1:50; SA 1:500'' (red, wide pink band) ≈5x10\textsuperscript{3} at ≈40 µm, ≈3x10\textsuperscript{3} at 100 µm, ≈2.5x10\textsuperscript{3} at 200--400 µm, ≈2x10\textsuperscript{3} at 500--700 µm. (c) y 0 to 4x10\textsuperscript{4}: ``SHORT NeuN 1:50; SA 1:200'' (black) ≈2.5x10\textsuperscript{4} at ≈25 µm, ≈1x10\textsuperscript{4} at 100 µm, ≈1x10\textsuperscript{3} by 500 µm; ``SHORT NeuN 1:50; SA 1:500'' (red) ≈5x10\textsuperscript{3} at ≈30 µm then ≈2x10\textsuperscript{3}; ``SHORT NeuN 1:50; SA 1:1000'' (blue, cyan band) ≈3.3x10\textsuperscript{4} at ≈25 µm, ≈1x10\textsuperscript{4} at 80 µm, near zero beyond 250 µm. (d) y 0 to 4x10\textsuperscript{4}: ``SHORT NeuN 1:100; SA 1:200'' (black) ≈8x10\textsuperscript{3} at ≈40 µm decaying slowly to ≈1x10\textsuperscript{3} by 500 µm; ``SHORT NeuN 1:100; SA 1:500'' (red) ≈4x10\textsuperscript{3} at ≈30 µm then ≈2x10\textsuperscript{3} decaying; ``SHORT NeuN 1:100; SA 1:1000'' (blue) ≈1x10\textsuperscript{3} flat with ≈1.5x10\textsuperscript{3} at ≈650 µm. (e) y 0 to 4x10\textsuperscript{4}: ``SHIELD NeuN 1:50; SA 1:200'' (black) ≈2.8x10\textsuperscript{4} at ≈25 µm, ≈1x10\textsuperscript{4} at 100 µm, ≈5x10\textsuperscript{3} at 200 µm, ≈1x10\textsuperscript{3} beyond 500 µm; ``SHIELD NeuN 1:50; SA 1:500'' (red) ≈4x10\textsuperscript{3} at ≈30 µm then ≈2x10\textsuperscript{3}; ``SHIELD NeuN 1:50; SA 1:1000'' (blue) ≈3x10\textsuperscript{3} at ≈20 µm, ≈1x10\textsuperscript{3}, rising to ≈4x10\textsuperscript{3} at ≈630 µm (the far edge). (f) y 0 to 4x10\textsuperscript{4}: ``SHIELD NeuN 1:100; SA 1:200'' (black) ≈3x10\textsuperscript{3} nearly flat; ``SHIELD NeuN 1:100; SA 1:500'' (red) ≈3x10\textsuperscript{3}, slightly decaying; ``SHIELD NeuN 1:100; SA 1:1000'' (blue, cyan band) ≈1.3x10\textsuperscript{4} at ≈50 µm decaying to ≈2x10\textsuperscript{3} by 200 µm.}
\figuredata{b}{figures/sfig3b.csv}{figures/sfig3b.py}{vector paths in the PDF (tier 2), exact up to the PDF's precision: one tiny marker per data point gives the mean, the two polylines of the series colour give the ±SD band edges, axes calibrated on the labelled tick marks (residuals, band-symmetry check and uncertainty in the CSV header)}
\figuredata{c}{figures/sfig3c.csv}{figures/sfig3c.py}{vector paths in the PDF (tier 2), exact up to the PDF's precision: one tiny marker per data point gives the mean, the two polylines of the series colour give the ±SD band edges, axes calibrated on the labelled tick marks (residuals, band-symmetry check and uncertainty in the CSV header)}
\figuredata{d}{figures/sfig3d.csv}{figures/sfig3d.py}{vector paths in the PDF (tier 2), exact up to the PDF's precision: one tiny marker per data point gives the mean, the two polylines of the series colour give the ±SD band edges, axes calibrated on the labelled tick marks (residuals, band-symmetry check and uncertainty in the CSV header)}
\figuredata{e}{figures/sfig3e.csv}{figures/sfig3e.py}{vector paths in the PDF (tier 2), exact up to the PDF's precision: one tiny marker per data point gives the mean, the two polylines of the series colour give the ±SD band edges, axes calibrated on the labelled tick marks (residuals, band-symmetry check and uncertainty in the CSV header)}
\figuredata{f}{figures/sfig3f.csv}{figures/sfig3f.py}{vector paths in the PDF (tier 2), exact up to the PDF's precision: one tiny marker per data point gives the mean, the two polylines of the series colour give the ±SD band edges, axes calibrated on the labelled tick marks (residuals, band-symmetry check and uncertainty in the CSV header)}

\pdfpage{5}
\begin{center}
\includegraphics[width=0.95\linewidth,height=0.6\textheight,keepaspectratio]{figures/sfig4.png}
\end{center}
\noindent \textbf{Supplementary Figure 4. Comparison of autofluorescence signal between SHIELD, SWITCH and SHORT.} SWITCH-, SHORT-, and SHIELD-processed slices excited at 405, 488, 561 and 638 nm. SHIELD and SHORT show similar autofluorescence signal at 405, 561, and 638 nm, while SHORT shows higher intensity at 488 nm (see also Fig. 1A). Laser power, 5 mW. The LUTs are fixed for each treatment. Scale bar = 100 μm.
\alttext{Supplementary Figure 4 (p.\,5): a grid of twelve LSFM strips, columns headed ``SHIELD'', ``SWITCH'', ``SHORT'' and rows labelled 405, 488, 561, 638 (vertical labels at the left), each strip a cross-section of the slice with the illuminated edge at the left, in blue, green, yellow and red look-up tables respectively; white scale bar (100 µm) in the top-left strip. At the right, under the header ``Merge channels'', three overlay images labelled ``SHIELD'', ``SWITCH'', ``SHORT''. The 405 row is almost black with a faint blue left edge in all three; in the 488 row the SWITCH strip is saturated bright green over its left half while SHIELD and SHORT are dim green with scattered puncta; in the 561 row SWITCH is bright yellow at the left, SHIELD and SHORT show yellow puncta on a dark background; in the 638 row SWITCH is brighter red with many puncta, SHIELD and SHORT show sparse red puncta. The SWITCH merge is bright yellow, the SHIELD and SHORT merges green with yellow/orange puncta. No quantitative axes.}

\pdfpage{6}
\begin{center}
\includegraphics[width=0.95\linewidth,height=0.6\textheight,keepaspectratio]{figures/sfig5.png}
\end{center}
\noindent \textbf{Supplementary Figure 5. High-resolution images acquired by LSFM (resolution of 0.55 x 0.55 x 3.3 μm).} Representative images showing morphological details of calretinin (CR)-, glutamic acid decarboxylase 67 (GAD67)-, vasoactive intestinal peptide (VIP)-, somatostatin (SST)-, calbindin (CB)-, parvalbumin (PV), neuronal nuclear protein (NeuN)-, glial fibrillary acidic protein (GFAP)-immunoreactive cortical neurons, and vimentin (VIM) immunolabeling of blood vessels. Scale bar = 100 μm for all images, except for GFAP (scale bar 10 = μm).
\alttext{Supplementary Figure 5 (p.\,6): nine greyscale LSFM micrographs. Eight square images in two rows of four, each labelled below: CR (bright somata with thin processes, horizontal streak artefacts), GAD67 (bright round somata and puncta on black), VIP (small bright cells and vertical processes), SST (a few bright cells with dendrites), CB (scattered bright cells), PV (bright cells with a web of thin processes), NeuN (many bright pyramidal-shaped somata with apical dendrites, tilted), Vim (bright tubular vessel profiles); white scale bar (100 µm) in the CR image. At the right a larger image labelled ``GFAP'': one star-shaped astrocyte with many radiating processes filling the field; white scale bar at its top left (10 µm). No quantitative axes.}

\pdfpage{7}
\begin{center}
\includegraphics[width=0.95\linewidth,height=0.6\textheight,keepaspectratio]{figures/sfig6.png}
\end{center}
\noindent \textbf{Supplementary Figure 6. Propidium iodide (PI) allows an efficient nuclear labeling in SHORT-processed samples.} Human motor cortex labeled with PI and acquired using LSFM. MIP of 20 slices. Excitation light 561 nm. Scale bar = 4 mm. Magnified image of the nuclear staining, scale bar = 200 μm.
\alttext{Supplementary Figure 6 (p.\,7): two red fluorescence images. Left: a whole slice of motor cortex (a gyrus seen in cross-section, with a sulcus entering from below and a dark notch at its tip), nuclei as a fine red speckle over the whole section, brighter blood vessels as elongated streaks, a white dashed square near the bottom of the sulcus marking the magnified region, and a white scale bar at the lower right (4 mm). Right: the magnified region, densely packed small red nuclei and a few vessel profiles; white scale bar (200 µm). No quantitative axes.}

\pdfpage{8}
\begin{center}
\includegraphics[width=0.95\linewidth,height=0.6\textheight,keepaspectratio]{figures/sfig7.png}
\end{center}
\noindent \textbf{Supplementary Figure 7. DiD staining on a SHORT-processed human brain slice.} Downscaled images of the precentral cortex (100 μm-thick) and the intensity signal of the fiber stained with DiD. SHORT is compatible with lipophilic dye and the signal along the thickness is sufficient to detect the labeled fibers. Excitation light, 638 nm; laser power, 1 mW. Scale bar = 10 μm.
\alttext{Supplementary Figure 7 (p.\,8): greyscale LSFM images. Left, a large image titled ``DiD staining – LSFM image'' (white text at the top left): a wedge-shaped piece of cortex, brighter (white) in its lower, deeper part and grey in the upper part, with many small dark holes (vessel profiles) and fine fibre texture; white scale bar labelled ``500 µm'' at the lower left. Right, three square panels stacked vertically, labelled ``0 µm'', ``50 µm'' and ``100 µm'' (depth): at 0 µm a bright field with many dark puncta and a few dark vessel profiles and a white scale bar (10 µm); at 50 µm a greyer field with fibre-like streaks and a bright cell; at 100 µm a dim, low-contrast field with faint fibres. No quantitative axes.}

\pdfpage{9}
\begin{center}
\includegraphics[width=0.95\linewidth,height=0.6\textheight,keepaspectratio]{figures/sfig8.png}
\end{center}
\begin{center}\textbf{Costaining NeuN (green) – Gad67 (red)}\end{center}
\decompnote{Title printed below the two images, above the legend.}

\noindent \textbf{Supplementary Figure 8. Comparison between the classic SWITCH protocol and SHORT for NeuN (green) - GAD67 (red) costaining.} The white arrow shows the high-resolution imaging of a typical costained neuron positive for NeuN and GAD67 in SHORT-processed slices. Scale bar = 100 μm. Excitation light: 488 nm and 638 nm.
\alttext{Supplementary Figure 8 (p.\,9): two square LSFM images, titled ``Classic SWITCH processed slice'' (left) and ``SHORT processed slice'' (right), with the printed title ``Costaining NeuN (green) – Gad67 (red)'' below them. Left: a green field with bright green NeuN somata, horizontal streak artefacts and almost no red signal; white scale bar (100 µm) at the lower right. Right: green somata on a green background with many red GAD67 puncta and several orange/yellow double-labelled cells; a white arrow near the upper left points at one orange (NeuN + GAD67) cell body. No quantitative axes.}

\pdfpage{10}
\begin{center}
\includegraphics[width=0.95\linewidth,height=0.6\textheight,keepaspectratio]{figures/sfig9.png}
\end{center}
\noindent\textbf{Stripping efficiency}
\decompnote{Title printed inside the figure, above the bar chart.}

\noindent \textbf{Supplementary Figure 9. Stripping efficiency in SHORT-processed human slices.} Tissue sections were processed with SHORT and were acquired using a Nikon C2 laser-scanning confocal microscope. Then, the samples were immunostained for NeuN with AlexaFluor 568 and were re-acquired by the same microscope. The samples were then stripped using the elution buffer (see Methods) for 4 °h at 80 °C and were acquired again. Quantification of fluorescence intensity seen in n = 3 samples. Error bars: standard deviation. Scale bar = 50 μm.
\alttext{Supplementary Figure 9 (p.\,10): a vector bar chart titled ``Stripping efficiency'' (printed inside the plot, upper right) and three raster images. The chart: y axis ``Intensity'' 0,0 to 3,5x10\textsuperscript{3} (ticks every 5,0x10\textsuperscript{2}, decimal commas), three open bars labelled ``No label/strip'', ``Anti-NeuN staining'', ``Post strip'', each with a SD error bar (upper half) and 40 grey filled circles (individual values; the legend says n = 3 samples) scattered over the bar. Approximate values: No label/strip ≈1.48x10\textsuperscript{3} (SD ≈0.22x10\textsuperscript{3}; points 1.2--1.8x10\textsuperscript{3}), Anti-NeuN staining ≈2.52x10\textsuperscript{3} (SD ≈0.7x10\textsuperscript{3}; points 1.8--3.5x10\textsuperscript{3}), Post strip ≈0.83x10\textsuperscript{3} (SD ≈0.22x10\textsuperscript{3}; points 0.5--1.15x10\textsuperscript{3}). Below the chart three red confocal images, one under each bar: the unlabelled slice (dim red with faint cell outlines), the NeuN-stained slice (bright red pyramidal somata), and the stripped slice (dim, cells barely visible); white scale bar (50 µm) at the bottom of the first image.}
\figuredata{(bar chart)}{figures/sfig9.csv}{figures/sfig9.py}{vector paths in the PDF (tier 2: bar outlines, whiskers, caps and circle fills), exact up to the PDF's precision; y axis calibrated on the eight labelled ticks (residual 0.46 units); 40 points per bar are drawn although the legend says n = 3 samples}

\pdfpage{11}
\begin{center}
\includegraphics[width=0.95\linewidth,height=0.6\textheight,keepaspectratio]{figures/sfig10.png}
\end{center}
\noindent \textbf{Supplementary Figure 10. Distortion analysis.} SHORT-processed human cortical slice immunostained for SST after the clearing process (post-clearing, red images). Then, the antibodies were stripped and restained for SST (stripping, green image). \textbf{a} Comparison of the structural similarity between two high-resolution images of SST-immunoreactive neurons after clearing and stripping process using the structural similarity index measure (SSIM index; range value between -1 and 1; two identical images have an SSIM index close to 1). The data were shown as the mean ± SD (N = 3). \textbf{b} Comparison of tortuosity of SST branches between high-resolution images post-clearing (red) and after stripping (red). The minimum distance (D) and the real branch length (d) between the two ends of each individual branch were measured to evaluate the morphological changes using the torsional index (d/D). The branch alteration after the stripping process is not statically significant (N.S.; N = 6). LSFM images; scale bar = 10 μm; scale bar of the merged image of SST branch = 5 μm.
\alttext{Supplementary Figure 10 (p.\,11): panels a and b, each with raster images at the left and a vector plot at the right. (a) Three square images of the same SST-immunoreactive neuron, headed ``Post-clearing'' (red), ``Stripping'' (green) and ``Merged'' (the overlay, yellow where both coincide: the soma and main dendrites coincide closely); white scale bar (10 µm) in the first. Plot: y axis ``SSIM index'' 0,0--1,0 (ticks every 0,2), one open bar labelled ``SST neurons'' with its top at ≈0.90, a SD error bar reaching ≈0.96, and three red filled circles at ≈0.88, ≈0.96 and ≈0.85. (b) A red and a green image of another neuron, each with a white dashed rectangle around one dendritic branch; a narrow merged image of that branch (red and green segments interleaved along the same path, with a red blob at the bottom); white scale bars (10 µm in the red/green images, 5 µm in the merged image). Beside them a sketch: two blue dots joined by a wavy blue line labelled ``d'' (the real branch length, in blue) and a straight red line labelled ``D'' (the minimum distance, in red). Plot: y axis ``Tortuosity'' 1,00--1,25 (ticks every 0,05), two x categories ``Post-clearing'' (grey squares) and ``Stripping'' (red squares), six pairs joined by black lines; a bracket labelled ``N.S.'' spans the two categories at the top. The pairs, approximately (post-clearing → stripping): 1.18 → 1.11, 1.09 → 1.23, 1.08 → 1.07, 1.07 → 1.06, 1.04 → 1.04, 1.02 → 1.02.}
\figuredata{a}{figures/sfig10a.csv}{figures/sfig10a.py}{vector paths in the PDF (tier 2: bar rectangle, error-bar strokes and the three red marker fills), exact up to the PDF's precision; y axis calibrated on the six labelled ticks (residual in the CSV header)}
\figuredata{b}{figures/sfig10b.csv}{figures/sfig10b.py}{vector paths in the PDF (tier 2: square marker centres, pairs from the connecting line segments), exact up to the PDF's precision; y axis calibrated on the six labelled ticks}

\pdfpage{12}
\begin{center}
\includegraphics[width=0.95\linewidth,height=0.6\textheight,keepaspectratio]{figures/sfig11.png}
\end{center}
\noindent \textbf{Supplementary Figure 11. Sample holder assembly. a} The sandwich holder consists of a cover glass, a spacer (thickness of 500 μm) and fused silica glass. The cleared specimen is fixed between the cover glass and the fused silica glass using a bicomponent glue. Such strategy allows soaking the sample in the TDE/PBS solution supplemented with a nuclear dye, and performing high-throughput imaging of several labelled human slices. Resliced images of mesoscopic reconstruction acquired by LSFM using the classic fused silica glass (\textbf{b}) upon the slice or the sandwich apparatus (\textbf{c}). Scale bar = 1 mm.
\alttext{Supplementary Figure 11 (p.\,12): (a) a schematic and nine photographs. Top left, the schematic: a blue outlined ``Coverglass'' slide with the slice (yellow outline) lying on it, a darker blue ``Fused quartz glass'' slide tilted above it with a curved arrow showing it being lowered, and below, after a black down-arrow, the closed sandwich with the slice between the two glasses. Top middle: a photograph of the parts on a white tray, labelled with arrows: two syringes (yellow and blue) ``Biocompatible bicomponent glue'', a ``Coverglass'', a rectangular frame ``Spacer'', and a white strip ``Fused quartz glass (RI 1.46)''. Top right: the assembled sandwich held by a green and a yellow binder clip (``Clip''), with the ``Cleared specimen'' (a yellowish slice) and a ``TDE drop'' on it, in a hexagonal dish. Middle row: three photographs of gloved hands applying the glue along the sandwich edge with a syringe (two views), and the closed sandwich with the slice inside. Bottom row: a syringe injecting ``TDE supplemented with nuclear staining'' into the sandwich, and two views of the filled sandwich with the slice. Far right, headed ``High-throughput imaging'': a photograph of a long rectangular holder on millimetre paper with three yellow, transparent slices mounted in a row. (b) A long horizontal greyscale resliced image (xz) of a slice imaged under a plain fused silica glass: the upper surface is wavy (undulating by roughly a slice thickness along the length); white scale bar at the top left (1 mm). (c) The same kind of resliced image with the sandwich: a flat, uniform band; white scale bar at the top right (1 mm). No quantitative axes.}

\pdfpage{13}
\begin{center}
\includegraphics[width=0.95\linewidth,height=0.6\textheight,keepaspectratio]{figures/sfig12.png}
\end{center}
\decompnote{The significance marks printed in the bar chart (``**'' over the first two green bars, ``*'' over the first two red bars) are text in the PDF; see the description.}
\noindent \textbf{Supplementary Figure 12. Signal retention (signal/background) after 1 year in the sandwich holder.} Quantification of the signal retention of 10 neurons acquired after immunostaining (Fig. 4l-n), and then re-acquired after 1 year, and 1 year+1 week kept in the sandwich holder. The data show a reduction of the fluorescence signal of 51\% for Alexa Fluor 488 (P < 0.01), 24\% for Alexa Fluor 568 (no significant), and 14\% for Alexa Fluor 647 (P < 0.05) (Mann-Whitney test). Scale bar xy = 1000 μm; scale bar xz = 500 μm; scale bar yz = 250 μm. The magnified image below the downscaled reconstruction highlights that after a year of storage the markers are well detectable. Scale bar = 100 μm.
\alttext{Supplementary Figure 12 (p.\,13): a vector bar chart (top left), three raster reconstructions (bottom left) and three raster close-ups (right). Chart: y axis ``Signal retained'' 0--8 (ticks every 2; axis to about 9), nine bars in three colour groups (green = CR, yellow = SST, red = NeuN), each group with the x labels ``Acquisition after staining'', ``Acquisition after 1 year'', ``LSFM bleaching'' (two-line labels), SD error bars (dark grey) and ten grey filled ellipses per bar (the individual neurons); significance bars with ``**'' over the first two green bars and ``*'' over the first two red bars. Approximate values (mean, SD): CR 5.1 (2.9; points from 2.6 to 8.1), 3.0 (1.2), 2.5 (0.4); SST 3.5 (1.3), 2.7 (1.5), 2.5 (1.4); NeuN 2.8 (0.5), 2.4 (0.4), 1.7 (0.3). Below the chart, three whole-slice MIP reconstructions of the same U-shaped frontal-cortex slice in green (CR), yellow (SST) and red (NeuN), each with a vertical yz strip at its right and a horizontal xz strip below, small axis glyphs (x/y, z/x, z/y) at the first one, and small white scale bars (1000 µm in xy, 500 µm in xz, 250 µm in yz per the caption). Right column: three magnified images labelled ``CR'' (green puncta and cells), ``SST'' (yellow cells on a striped background), ``NeuN'' (dense red somata), with a white scale bar (100 µm) in the CR image.}
\figuredata{(bar chart)}{figures/sfig12.csv}{figures/sfig12.py}{vector paths in the PDF (tier 2: bar rectangles, whisker and cap strokes, ellipse marker fills), exact up to the PDF's precision; y axis calibrated on the five labelled ticks (residual in the CSV header)}

\pdfpage{14}
\begin{center}
\includegraphics[width=0.95\linewidth,height=0.6\textheight,keepaspectratio]{figures/sfig13.png}
\end{center}
\decompnote{The table printed at the top of the figure (Arial text in the PDF), transcribed here and in figures/sfig13.csv:}
\begin{center}\begin{tabular}{lcc}\toprule
Shrinking effect & PBS (cm\textsuperscript{2}) & SHORT (cm\textsuperscript{2})\\\midrule
Sample 1 & 7.188 & 7.048\\
Sample 2 & 7.181 & 6.92\\
Sample 3 & 7.396 & 7.027\\\bottomrule\end{tabular}\end{center}

\noindent \textbf{Supplementary Figure 13. Shrinking characterization of SHORT processed slices after incubation in the 68\% TDE/PBS solution (N = 3 samples).} The slices to be enclosed into the holder, were equilibrated in the TDE solution, tuned for the refractive index matching of the delipidated sample (RI ≈ 1.46 with TDE/PBS solution at 68\%), minimizing optical aberration and making the sample almost completely transparent. After equilibration of the processed sample into the PBS/TDE solution, we observed an isotropic tissue shrinkage of a factor of 3.5 ± 1.5 due to the RI matching obtained with the organic solvent TDE. Generally, the TDE solution is supplemented with nuclear staining.
\alttext{Supplementary Figure 13 (p.\,14): at the top a three-column table with a grey header row ``Shrinking effect'', ``PBS (cm\textsuperscript{2})'', ``SHORT (cm\textsuperscript{2})'' and rows Sample 1: 7.188, 7.048; Sample 2: 7.181, 6.92; Sample 3: 7.396, 7.027 (transcribed above and in figures/sfig13.csv; the areas after SHORT are 2--4\% smaller than in PBS, which does not match the ``factor of 3.5 ± 1.5'' in the legend). Below, two photographs of a slice on millimetre paper: left in PBS, a pale, opaque, cream-coloured square slab; right after SHORT/TDE, the same slab yellow and transparent, the grid lines visible through it. The grid squares are not labelled, so the photographs give no scale beyond the millimetre paper.}
\figuredata{(table)}{figures/sfig13.csv}{(none: a printed table, no plot)}{printed values in the PDF text layer, exact}

\pdfpage{15}
\begin{center}
\includegraphics[width=0.95\linewidth,height=0.6\textheight,keepaspectratio]{figures/sfig14.png}
\end{center}
\noindent \textbf{Supplementary Figure 14}. \textbf{Multiplexed staining of SHORT-processed slice and the homogenous staining along the thickness.} xy and zx MIP images of the superior frontal cortex stained for SST (Alexa Fluor 568; orange), VIP (Alexa Fluor 647; blue), and CR (Alexa Fluor 488; grey). MIPs of 150 slices (whole thickness). Scale bar xy images = 1 cm; scale bar xz = 500 μm; scale bar of the magnified insets = 100 μm.
\alttext{Supplementary Figure 14 (p.\,15): three rows of raster LSFM images, one per marker, each labelled at the top left of its first image: ``SST'' (orange look-up table), ``VIP'' (blue), ``CR'' (grey). In each row: a large xy maximum-intensity projection of the whole superior-frontal-cortex slice (an elongated slab with a bump at the top), a white scale bar along its top (1 cm), a small x/y axis glyph at the lower right and a small dotted square (cyan for SST, orange for VIP and CR) marking the inset position; under it a long thin xz strip (the slice seen edge-on, uniformly bright through its thickness) with a white scale bar at its left (500 µm), a z/x axis glyph at its right and a dotted rectangle at the inset position; at the right two magnified insets framed with dashed lines, labelled ``xy'' and ``xz'' in colour, with white scale bars (100 µm): bright cells and processes distributed through the whole depth of the xz inset. No quantitative axes.}

\pdfpage{16}
\begin{center}\small\begin{longtable}{llllll}\toprule
\textbf{Molecule} & \textbf{Company} & \textbf{Cat. n.} & \textbf{Host} & \textbf{P/M} & \textbf{Dilution}\\\midrule
NeuN & Merck & ABN91 & Chicken & P & 1:50\\
GAD67 & Abcam & ab26116 & Mouse & M & 1:200\\
PV & Abcam & ab11427 & Rabbit & P & 1:200\\
PV & Abcam & ab32895 & Goat & P & 1:200\\
CB & Abcam & ab207528 & Rabbit & M & 1:200\\
VIP & Abcam & ab214244 & Rabbit & M & 1:200\\
SST & Abcam & ab30788 & Rat & M & 1:200\\
NPY & Abcam & ab6173 & Sheep & P & 1:200\\
NPY & Abcam & ab11247 & Mouse & M & 1:200\\
SMI-32 & Merck & NE1023 & Mouse & M & 1:200\\
SMI-31 & Eurogentec OptimAb & SMI-31P-050 & Mouse & M & 1:300\\
Neurofilament & Abcam & ab4680 & Chicken & P & 1:200\\
GluS & Merck & MAB302 & Mouse & M & 1:200\\
MAP2 & Abcam & ab5392 & Chicken & P & 1:200\\
GFAP & Abcam & ab194324 & Rabbit & M & 1:200\\
Iba1 & Abcam & ab195031 & Rabbit & M & 1:200\\
Coll IV & Abcam & ab6586 & Rabbit & P & 1:200\\
Vim & Abcam & ab8069 & Mouse & M & 1:200\\
CR & Proteintech & 66496-1-Ig & Mouse & M & 1:200\\
CR & Proteintech & 12278-1-AP & Rabbit & P & 1:200\\
Anti-Rat IgG, AF 568 & Abcam & ab175475 & Donkey & P & 1:200\\
Anti-Rabbit IgG, AF 568 & Abcam & ab175470 & Donkey & P & 1:200\\
Anti-Chicken IgY, AF 568 & Abcam & ab175711 & Goat & P & 1:200\\
Anti-Mouse IgG, AF 568 & Abcam & ab175700 & Donkey & P & 1:200\\
Anti-Sheep IgG, AF 568 & Abcam & ab175712 & Donkey & P & 1:200\\
Anti-Rabbit IgG, AF 488 & Abcam & ab150077 & Goat & P & 1:200\\
Anti-Chicken IgY, AF 488 & Abcam & ab150169 & Goat & P & 1:200\\
Anti-Rabbit IgG, AF 488 & Jackson & 611-545-215 & AffiniPure Alpaca & P & 1:200\\
Anti-Goat IgG, AF 488 & Jackson & 805-545-180 & AffiniPure Bovine & P & 1:200\\
Anti Rabbit IgG, AF 647 & Jackson & 611-605-215 & AffiniPure Alpaca & P & 1:200\\
Anti Mouse IgG, AF 647 & Abcam & ab150107 & Donkey & P & 1:200\\
Anti Chicken IgY, AF 647 & Abcam & ab150171 & Goat & P & 1:200\\
DiD & Thermo Fisher Scientific & D7757 &  &  & 0.5 mg/ml\\
DAPI (Dilactate) & Thermo Fisher Scientific & D3571 &  &  & 1:50\\
SYTOX GREEN & Thermo Fisher Scientific & S7020 &  &  & 1:100\\
Propidium Iodide & Thermo Fisher Scientific & P3566 &  &  & 1:50\\
\bottomrule\end{longtable}\end{center}
\noindent \textbf{Supplementary Table 1. Antibodies compatible with SHORT and the used dilutions.}
\decompnote{Also tables/stable1.csv.}

\pdfpage{17}
\begin{center}\begin{tabular}{lllll}\toprule
 & \textbf{Area} & \textbf{Age} & \textbf{Fixative} & \textbf{Fixation}\\\midrule
\textbf{Donor 1} & - Precentral gyrus - Hippocampus & 99 years old & 10\% Formalin & 6 months\\
\textbf{Donor 2} & - Prefrontal cortex & 62 years old & 10\% Formalin & 4 years\\
\textbf{Donor 3} & - Broca’s area - Motor cortex & 79 years old & 10\% Formalin & 10 years\\
\bottomrule\end{tabular}\end{center}
\noindent \textbf{Supplementary Table 2. Area, age, fixative molecule, and fixation of the human brain material used in this work.}
\decompnote{Also tables/stable2.csv. The first column has no header in the PDF.}


\sourcefile{pesce2022\_supplementary\_file\_descriptions.pdf}
\decompnote{Separate one-page publisher file (42003\_2022\_3390\_MOESM3\_ESM.pdf) listing the supplementary movies. Supplementary Movies 1--4 themselves were not downloaded and are not in the source folder.}
\pdfpage{1}
\begin{center}\textbf{Description of Additional Supplementary Files}\end{center}

\noindent \textbf{File name:} Supplementary Movie 1 \\ \textbf{Description:} High resolution stacks of SHORT-processed slices acquired using LSFM (resolution of 0.55 x 0.55 x 3.3 μm) labelled for CR with Alexa Fluor 488

\noindent \textbf{File name:} Supplementary Movie 2 \\ \textbf{Description:} High resolution stacks of SHORT-processed slices acquired using LSFM (resolution of 0.55 x 0.55 x 3.3 μm) labelled for SST with Alexa Fluor 568

\noindent \textbf{File name:} Supplementary Movie 3 \\ \textbf{Description:} High resolution stacks of SHORT-processed slices acquired using LSFM (resolution of 0.55 x 0.55 x 3.3 μm) labelled for NeuN with Alexa Fluor 647.

\noindent \textbf{File name:} Supplementary Movie 4 \\ \textbf{Description:} Downscaled reconstruction of the Broca’s area labelled for NeuN with Alexa Fluor 647. Resolution of 3.3 x 3.3 x 3.3 μm.

\end{document}
